\section{The Mersenne branch and the arithmetic zoo}\label{sec:zoo}

Beyond the Clay problems the engine reading applies, lightly, to a whole zoo of classical
questions about primes. This section is a survey: for each family we record what is
genuinely proved (an embedding into the $6m\pm1$ language, a twin criterion, a classical
divisibility), what would a refutation cost (a manifestation front mirroring
\cref{sec:riemann}), and where the honesty stops. The unifying discipline is a deliberate
restraint: for none of these families is a decree boundary taken. Every trilemma is
cleared, every boundary is machine-admissible, and every one is left open on purpose ---
a choice we make plain below. Throughout, classical facts are those of Euclid's
\emph{Elements}~\cite{euclid-elements} and Hardy--Wright~\cite{hardy-wright}.

\subsection{The Mersenne branch: an honest bridge}

Mersenne numbers $2^p-1$ look like ready-made inhabitants of the $6m\pm1$ grid, and the
temptation is to declare that the twin engine solves them ``along the way''. It does not.
We record first the genuine embedding.

\begin{theorem}[Mersenne on the grid]\label{thm:mersenne_eq_sixCenter_add_one}
\green{}\; For odd $p$, writing $m_p=(2^{p-1}-1)/3$,
\begin{equation}\label{eq:mersenne-embed}
  2^p-1 = 6\,m_p+1.
\end{equation}
\leantag{mersenne\_eq\_sixCenter\_add\_one}
\end{theorem}

Thus every odd Mersenne number is exactly the plus side of the centre $m_p$; its lower side
$6m_p-1 = 2^p-3$ gives an immediate twin criterion.

\begin{theorem}[Mersenne twin criterion]\label{thm:isTwinCenter_mersenneCenter_iff}
\green{}\; For odd $p\ge 2$ the centre $m_p$ is a twin centre if and only if both $2^p-3$
and $2^p-1$ are prime. \leantag{isTwinCenter\_mersenneCenter\_iff}
\end{theorem}

Such Mersenne twins exist ($p=3$ gives $(5,7)$, $p=5$ gives $(29,31)$), and the only honest
arrow between the topics runs \emph{from} them: a subsequence of twins is still twins.

\begin{theorem}[unbounded Mersenne twins feed the twins]\label{thm:twinLowersInfinite_of_mersenneTwins}
\green{}\; If the Mersenne twin centres are unbounded, then the lower twin centres are
infinite. \leantag{twinLowersInfinite\_of\_mersenneTwins}
\end{theorem}

\emph{Proof idea.} The pairs $(2^p-3,\,2^p-1)$ are then unbounded twin pairs, and a
subsequence of twins is carried across the honest bridge to the programme's genuine goal.
The arrow is trivial and points from the stronger hypothesis to the weaker conclusion; the
reverse implication ``twins $\Rightarrow$ Mersenne'' is \emph{not} a theorem here --- the
branch carries an explicit coverage marker recording that no such implication is claimed, and
the goal ``Mersenne primes are infinite'' is \red{} open, derived from nowhere in the branch.

\begin{remark}
An earlier ``forward series'' of $34$ bricks aimed to prove Mersenne infinitude
unconditionally; an internal audit exposed it as \warn{}~vacuous --- its no-engine
packages were uninhabited, and its headline theorems followed from empty premises. It is
recorded, not retouched: the reduction chains here consume only genuine hypotheses.
\end{remark}

\subsection{A refutation presents an engine, and the boundary we did not take}

Run the Mersenne branch through the manifestation architecture of \cref{sec:riemann}. The
deviation is not a data point but a $\Pi$-witness of \emph{finiteness}: an absence witness
asserting that every prime pair $(2^p-3,\,2^p-1)$ sits with its lower side no higher than $P$.
The plumbing ties it to the goal exactly (unboundedness is equivalent to the absence of such a
witness, \green{}), and the witness cannot be forged: the centres form the base-$4$ chain
$m\to 4m+1$, whose first step $5\to 1$ carries no proper peel at any scale (\green{}), because
the leading coefficient $4$ is neither prime nor $\equiv\pm1\pmod 6$ --- sides do not map to
sides. So no unconditional supply of flows is buildable here, exactly the Riemann situation.

\begin{theorem}[a Mersenne refutation exhibits an engine]\label{thm:mersenneRefutation_carries_engine}
\green{}\; An absence witness, the Mersenne manifestation law, and reconciled books at a scale
no lower than $P$ present a perpetual engine as a concrete object, before any killing.
\leantag{mersenneRefutation\_carries\_engine}
\end{theorem}

\emph{Proof idea.} The resolving projection yields a stable, energy-free universe; the law
turns the absence into an infinite family of flows; a finite key is forced to collide two
of its members; from the collision an engine-witness is assembled. Composed with the no-engine
wall and the accepted boundary this forces the Mersenne twins unbounded, and across the
bridge the lower twin centres come out infinite (\green{}). The witness gate is load-bearing:
the ungated form at $P=0$ would produce a supply at a resolved scale, contradicting the green
impossibility \cref{thm:no_deviationFlowSupply_at_resolved_scale} of a supply there.

The trilemma is cleared and the boundary field \code{mersenneBoundary} is machine-admissible.
Its price is exact.

\begin{theorem}[exact price of the Mersenne decree]\label{thm:mersenneManifestationLaw_iff_unbounded_of_boundary}
\green{}\; Under the accepted boundary, the manifestation law is equivalent to the
unboundedness of the Mersenne twins:
\begin{equation}\label{eq:mersenne-price}
\begin{gathered}
\code{TheStrictLastStep00Obligation}\ \Rightarrow\\
\bigl(\code{MersenneManifestationLaw}\iff\code{MersenneTwinCentersUnbounded}\bigr).
\end{gathered}
\end{equation}
\leantag{mersenneManifestationLaw\_iff\_unbounded\_of\_boundary}
\end{theorem}

The decree would carry exactly Mersenne-twin strength --- and here, for the first time in the
programme, the sign of the heuristic is \emph{inverted}: the expected number of pairs
$(2^p-3,\,2^p-1)$ is a convergent sum, only $p=3,5$ are known, and for every $p\equiv3\pmod4$
the number $2^p-3$ is divisible by $5$. To take the boundary would be to bet against the
heuristic; we \emph{deliberately did not take it}. The machine criterion says ``you may'';
honesty says ``may is not should''.

\subsection{Grid sides and Polignac patterns: singletons green, pairs conjectural}

The pair problems --- twins (gap $2$), cousins (gap $4$), sexies (gap $6$), and de Polignac in
general --- decompose geometrically on the grid. First the spine of honesty: each side
separately carries infinitely many primes.

\begin{theorem}[each grid side carries infinitely many primes]\label{thm:minusSide_primes_unbounded}
\green{}\; Above any $n$ there is a prime $p>n$ with $p\bmod 6 = 5$; and, separately, a prime
$p>n$ with $p\bmod 6 = 1$. \leantag{minusSide\_primes\_unbounded}
\end{theorem}

\emph{Proof idea.} This is Dirichlet's theorem on primes in arithmetic progressions, borrowed
whole from \code{mathlib}~\cite{mathlib}: the classes $1$ and $5$ modulo $6$ are coprime to
$6$, so each carries infinitely many primes. About one and the same centre the two columns say
nothing --- a marker (\green{}, an \code{abbrev := True}) signs this honestly, and the
coincidence of the sides at a shared centre is the twin conjecture, neither asserted nor
derived here.

Cousins and sexies are read off the same grid. A cousin pair $(p,p+4)$ crosses two adjacent
centres, $(6m+1,\,6(m+1)-1)$; a sexy pair $(p,p+6)$ stays on one side of two adjacent centres.
The geometry forces a residue.

\begin{theorem}[cousins are born on the plus side]\label{thm:cousin_mod_six}
\green{}\; Every junior member of a cousin pair with $p>3$ lies on the plus side:
$p\bmod 6 = 1$. \leantag{cousin\_mod\_six}
\end{theorem}

\emph{Proof idea.} A prime $p>3$ has $p\bmod 6\in\{1,5\}$; the case $5$ forces $p+4\equiv 0
\pmod 3$, impossible for the prime $p+4>3$. Cousins are born only on the plus side.

\begin{theorem}[adjacent twin centres yield a constellation]\label{thm:constellation_of_adjacent_twinCenters}
\green{} (even without \code{Classical.choice})\; Two adjacent twin centres $m$ and $m+1$
simultaneously yield the cousins on $m$ and both sexy pairs on $m$.
\leantag{constellation\_of\_adjacent\_twinCenters}
\end{theorem}

\emph{Proof idea.} Then all four sides $6m\mp1$, $6(m+1)\mp1$ are prime, and cousins
$(6m+1,\,6m+5)$, the minus-sexy $(6m-1,\,6m+5)$ and the plus-sexy $(6m+1,\,6m+7)$ are pure
rearrangements --- at $m=1$ this is the constellation $5,7,11,13$. The proof does not
\emph{construct} adjacent twin centres; that they exist arbitrarily high is a conjecture
stronger than twins. Both families run through the same manifestation front, both boundary
fields \code{cousinBoundary}, \code{sexyBoundary} are machine-admissible, and both are not
taken --- here there is no chain to saw at all, and a serial decree would erase the uniqueness
in which the quarantine's meaning resides.

\subsection{The plus/minus axis is orthogonal to proper time}

The two sides of a centre --- the would-be members of a pair $6m\pm1$ --- are exchanged by the side
flip \code{mirrorSide}, an involution (\code{mirrorSide\_involutive}). Proper time does not see it.

\begin{theorem}[the plus/minus flip is orthogonal to proper time]\label{thm:lexRank_side_invariant}
\green{}\ \code{lexRank} is invariant under the side flip: the two members of a pair $6m\pm1$ sit at
the same proper time, differing only along the orthogonal $\pm$ axis --- the exact complement of the
arrow of time (\cref{thm:timeArrow_step}). \leantag{lexRank\_side\_invariant}
\end{theorem}

\emph{Proof idea.} \code{lexRank} is $3c+p$, where the centre $c$ (\code{centerOf}) and the phase $p$
(\code{phase}) both discard the \code{Side} of a defect; the flip changes nothing the clock reads
(by \code{rfl}).

\begin{theorem}[the pair is a reflection about $6m$]\label{thm:signedSideValue_reflection_sum}
\green{}\ In $\mathbb{Z}$, $(6z+1)+(6z-1)=12z$ exactly (\code{signedSideValue\_reflection\_sum}); the
$\mathbb{N}$ form holds for $m\ge 1$ (\code{sideValue\_reflection\_sum}), the guard avoiding the
truncation $6\cdot 0-1=0$ (the grave of zero). \leantag{signedSideValue\_reflection\_sum}
\end{theorem}

\begin{remark}[charge conjugation $C$, not time reversal $T$]
The $\pm$ side flip is a fixed-time involution ($C^2=1$), a discrete analogue of charge conjugation
$C$ --- \emph{not} a time reversal $T$. Reading the pair as opposite \emph{time} directions would be a
\code{lexRank} increase, i.e.\ the forbidden perpetual engine (the arrow of time,
\cref{thm:timeArrow_step}); the orthogonality above is exactly why the $C$ reading is the consistent
one --- an internal symmetry orthogonal to the temporal axis in the sense of Coleman--Mandula. This is
a translation, not a claim: the Lean has no Hilbert space, field or charge, only a two-element
\code{Side}.
\end{remark}

\subsection{Sophie Germain: doubling and a classical divisibility}

A Sophie Germain pair is $(p,\,2p+1)$, both prime. In grid coordinates it is one-sided, and
the SG map of centres is pure doubling.

\begin{theorem}[Sophie Germain centres double]\label{thm:sg_center_doubling}
\green{}\; $2(6m-1)+1 = 6(2m)-1$: an SG prime $p=6m-1$ has its companion $2p+1$ on the minus
side of centre $2m$. \leantag{sg\_center\_doubling}
\end{theorem}

The map $m\to 2m$ is the third sibling of the chains, beside Mersenne's $4m+1$ and the pentadic
$5m+1$. But the branch carries the one substantive classical theorem of the whole zoo,
machine-proved to the last step.

\begin{theorem}[a Sophie Germain companion divides its Mersenne number]\label{thm:sophieGermain_divides_mersenne}
\green{}\; If $p$ is prime, $p\equiv 3\pmod 4$, $p>3$, and $q=2p+1$ is prime, then $q\mid M_p$
and the Mersenne number $M_p=2^p-1$ is composite.
\leantag{sophieGermain\_divides\_mersenne}
\end{theorem}

\begin{proof}
This is the classical Euler--Lagrange route~\cite{hardy-wright}. Since $p\equiv3\pmod4$ we have
$q=2p+1\equiv7\pmod8$, which is exactly the condition $q\equiv\pm1\pmod8$ under which $2$ is a
quadratic residue modulo $q$ (\code{ZMod.exists\_sq\_eq\_two\_iff}). By Euler's criterion
$2^{(q-1)/2}\equiv 1\pmod q$, and $(q-1)/2=p$, so $2^p\equiv 1\pmod q$, i.e.\ $q\mid 2^p-1=M_p$.
For $p\ge4$ one has $2p+1<2^p-1$, so $1<q<M_p$ and
$M_p$ is composite. The proviso $p>3$ is honest: at $p=3$ the companion $q=7=M_3$ is itself the
Mersenne number and is prime, so the divisor coincides with the dividend.
\end{proof}

\begin{theorem}[a worked composite Mersenne number]\label{thm:mersenne_composite_examples}
\green{}\; $23\mid M_{11}=2047$, and $M_{11}$ is composite.
\leantag{mersenne\_composite\_examples}
\end{theorem}

Here $p=11\equiv3\pmod4$, $q=23$ is prime, and $2047=23\cdot89$ splits into exactly the predicted
divisor. Through the restricted $3\pmod4$ manifestation front this yields, conditionally on the
open Sophie Germain conjecture, unboundedly many composite Mersenne numbers (\green{},
conditional) --- a one-way bridge onto the \emph{composite} side of the same $M_p$ whose
\emph{prime} side the Mersenne boundary declined to bet on.

\subsection{Goldbach and Legendre: a decidable witness}

For Goldbach and Legendre the deviation is a data object that is moreover \emph{pointwise
decidable} --- the strongest form of unpresentability in the series. Each candidate is the
easiest thing to check, yet presenting a genuine witness is impossible.

\begin{theorem}[a Goldbach violation must sit at $E\ge52$]\label{thm:goldbachViolation_ge_52}
\green{}\; Every witness of a Goldbach violation sits at $E\ge 52$.
\leantag{goldbachViolation\_ge\_52}
\end{theorem}

\emph{Proof idea.} A Goldbach violation ``$E\ge4$ even, not a sum of two primes'' is decidable
through a bounded search; \code{decide} verifies that every even $E<52$ decomposes, so any fake
below $52$ dies on the spot. The literature bound $4\cdot10^{18}$ is not formalised; only
$\ge52$ is green. The manifestation law is object-quantified (no gate: the scale anchor
$M_0:=E$ is the number itself), and the essence carries through.

\begin{theorem}[Goldbach from manifestation and boundary]\label{thm:noGoldbachViolation_of_manifestation_and_boundary}
\green{}\; No engines, an accepted boundary, and the manifestation law imply there are no
Goldbach violations --- the Goldbach conjecture holds.
\leantag{noGoldbachViolation\_of\_manifestation\_and\_boundary}
\end{theorem}

\emph{Proof idea.} A hypothetical witness $V$ yields the scale $M_0:=V.1$; the boundary supplies
a resolving projection there; the law turns the violation into an infinite family of flows, from
whose collision an engine-witness is assembled --- killed by ``no engines''. Across the bridge
tying violations to the conjecture, this reaches the conjecture itself. Legendre's front
is identical (\green{}), but carries a unique unconditional companion --- Bertrand's postulate.

\begin{theorem}[no prime desert survives a doubling]\label{thm:no_desert_doubles}
\green{}\; A prime desert cannot survive a doubling: for $n\neq0$ there is always a prime
between $n$ and $2n$. \leantag{no\_desert\_doubles}
\end{theorem}

\emph{Proof idea.} A direct repackaging of Bertrand's postulate
(\code{Nat.exists\_prime\_lt\_and\_le\_two\_mul}, \code{mathlib}~\cite{mathlib}, unconditional).
But it does \emph{not} solve Legendre, and we mark this machine-checkably.

\begin{theorem}[Legendre's interval is shorter than Bertrand's]\label{thm:legendre_interval_shorter_than_bertrand}
\green{}\; For $n\ge3$ one has $(n+1)^2<2n^2$: Legendre's interval $(n^2,(n+1)^2)$ is
\emph{shorter} than Bertrand's doubling $[n,2n]$.
\leantag{legendre\_interval\_shorter\_than\_bertrand}
\end{theorem}

Hence the green Bertrand is powerless on the quadratic crack, and the module records that no
Bertrand-to-Legendre implication is claimed. A green theorem next to an open problem
does not make it green; it delineates exactly where the proven ends.

\subsection{Perfect numbers: green Euclid--Euler both ways}

Both directions of the oldest theorem of the \emph{Elements}~\cite{euclid-elements} are now green,
re-exported from the \code{mathlib} archive~\cite{mathlib,wiedijk100}.

\begin{theorem}[Euclid's direction: Mersenne prime gives a perfect number]\label{thm:perfect_of_mersennePrime}
\green{}\; For $2\le p$ with $2^p-1$ prime, the number $2^{p-1}(2^p-1)$ is perfect --- Euclid's
direction (\emph{Elements} IX.36). \leantag{perfect\_of\_mersennePrime'}
\end{theorem}

\begin{theorem}[Euler's converse: every even perfect number is Euclid's]\label{thm:evenPerfect_eq}
\green{}\; Every even perfect number has the form $2^k(2^{k+1}-1)$ with a prime Mersenne factor
--- Euler's converse. \leantag{evenPerfect\_eq}
\end{theorem}

Together these two arrows are the entire classical theorem: Mersenne primes and even perfect
numbers are one object. This makes an ancient equivalence green.

\begin{theorem}[Mersenne primes and even perfect numbers rise together]\label{thm:mersennePrimesInfinite_iff_evenPerfectUnbounded}
\green{}\; The infinitude of the Mersenne primes is equivalent to the unboundedness of even
perfect numbers. \leantag{mersennePrimesInfinite\_iff\_evenPerfectUnbounded}
\end{theorem}

\emph{Proof idea.} Forward, Euclid's construction with $N<p<2^p\le 2^{p-1}(2^p-1)$ drives the
perfect numbers upward without a ceiling; backward, Euler's classification dissects an unbounded
even perfect number and a divisibility lemma forces a growing prime exponent.
Both sides are open; what is green is the \emph{equivalence} --- the oldest question of the
\emph{Elements} and the Mersenne question are literally one. The even side is constructive and
locked, so a deviation --- an odd perfect number --- can live only on the odd side; its witness is
pointwise decidable and $\ge101$ is green (the literature bound $>10^{2200}$ is not formalised).
Euler's classical form of the hypothetical witness is machine-verified.

\begin{theorem}[Euler's form of an odd perfect number]\label{thm:odd_perfect_euler_form}
\green{}\; Every odd perfect $n$ has the form $n=q^{\alpha}m^2$ with $q$ prime, $q\equiv1\pmod4$,
$\alpha\equiv1\pmod4$, and $q\nmid m$. \leantag{odd\_perfect\_euler\_form}
\end{theorem}

\emph{Proof idea.} For odd $n$ the factor $2$ enters $\sigma(n)=2n$ in exactly the first power;
$\sigma$ being multiplicative, exactly one prime power $\sigma(q^\alpha)$ can be even, forcing a
unique odd exponent, and mod-$4$ counting pins $q\equiv\alpha\equiv1\pmod4$. This is Euler's
eighteenth-century form; together with ``at least three distinct prime divisors'' it tightens the
domain of the nonexistent witness --- without asserting its existence either way.

\subsection{Fermat numbers: a second bet against expectations}

The Fermat numbers $F_k=2^{2^k}+1$ mirror Mersenne on the other side of the grid.

\begin{theorem}[a Fermat number is the minus side of its centre]\label{thm:six_mul_fermatCenter}
\green{}\; For $k\ge1$, with $c_k=(F_k+1)/6$, one has $6c_k=F_k+1$, i.e.\ $F_k=6c_k-1$: a Fermat
number is the \emph{minus} side of its centre. \leantag{six\_mul\_fermatCenter}
\end{theorem}

\emph{Proof idea.} For $k\ge1$ the exponent $2^k$ is even, so $F_k\equiv2\pmod3$ and, by oddness,
$F_k\equiv5\pmod6$; $F_0=3$ is off the grid and excluded. The chain of centres is here
\emph{quadratic}, $c_{k+1}=6c_k^2-4c_k+1$ (\green{}), which again cannot be sawed --- the step
$3\to1$ carries no proper peel (\green{}) --- so no forged witness exists. The
localisation of the absence witness is the record of the programme.

\begin{theorem}[the Fermat-twin absence bound is at least $65537$]\label{thm:fermatAbsenceBound_ge_65537}
\green{}\; Every Fermat-twin absence bound is at least $65537$.
\leantag{fermatAbsenceBound\_ge\_65537}
\end{theorem}

\emph{Proof idea.} The pair at $k=4$, namely $F_4=65537$ and $F_4+2=65539$, both prime, exists
greenly, and no witness can cut it off from below; note $k=3$ gives no pair, $F_3+2=259=7\cdot37$
composite. The rest of the front mirrors Mersenne word for word, and the price is exact.

\begin{theorem}[exact price of the Fermat decree]\label{thm:fermatManifestationLaw_iff_unbounded_of_boundary}
\green{}\; Under an accepted boundary the manifestation law is equivalent to the unboundedness of
Fermat twins. \leantag{fermatManifestationLaw\_iff\_unbounded\_of\_boundary}
\end{theorem}

The decree would carry exactly Fermat-twin strength, but the sign of the heuristic is inverted
harder than Mersenne's: the sum $\sum 2^{-2^k}$ converges almost instantly, only $F_0$ through
$F_4$ are known prime while $F_5$ through $F_{32}$ are provably composite, and the consensus expects
no new Fermat primes at all. Fermat himself bet that all $F_k$ are prime and lost to Euler's
$641\mid F_5$. So the field \code{fermatBoundary} is machine-admissible and, once more, deliberately
not taken. What for Mersenne was a one-off refusal by the sign of the heuristic becomes with Fermat
a \emph{rule} of the zoo: an honest boundary demands a greenly-unpresentable witness \emph{and} a
bet one need not be ashamed of.

\begin{remark}
None of the sections above declares any open problem solved. \code{twin\_prime\_conjecture} remains
\code{sorry}; Goldbach, Legendre, Sophie Germain, Polignac (gaps $4,6$), the infinitude of Mersenne
primes and of Fermat primes, and the existence of an odd perfect number all remain \red{} open; the
quarantine's taint is unchanged. What is green is compressed classical arithmetic and the honest
delineation of where each wager begins.
\end{remark}
